\documentclass[11pt,a4paper]{article}

\usepackage[utf8]{inputenc}
\usepackage[T1]{fontenc}
\usepackage{amsmath,amssymb,amsthm}
\usepackage{mathtools}
\usepackage{geometry}
\usepackage{hyperref}
\usepackage{booktabs}
\usepackage{longtable}
\usepackage{enumitem}
\usepackage{xcolor}
\usepackage{listings}
\usepackage{caption}
\usepackage{subcaption}
\usepackage{fancyhdr}

\geometry{margin=2.5cm}
\pagestyle{fancy}
\fancyhead[L]{\small Fourth-Derivative Gravity Propagators}
\fancyhead[R]{\small Adversarial Verification Report}

\hypersetup{
  colorlinks=true,
  linkcolor=blue!70!black,
  citecolor=green!50!black,
  urlcolor=blue!60!black
}

\lstset{
  basicstyle=\ttfamily\footnotesize,
  backgroundcolor=\color{gray!10},
  frame=single,
  breaklines=true,
  numbers=left,
  numberstyle=\tiny\color{gray},
  keywordstyle=\color{blue},
  commentstyle=\color{green!50!black},
  stringstyle=\color{red!60!black},
  language=Python
}

\newtheorem{theorem}{Theorem}[section]
\newtheorem{lemma}[theorem]{Lemma}
\newtheorem{claim}[theorem]{Claim}
\newtheorem{definition}[theorem]{Definition}
\newtheorem{remark}[theorem]{Remark}

\newcommand{\p}{\partial}
\newcommand{\lap}{\nabla^2}
\newcommand{\dbox}{\Box}
\newcommand{\TT}{\mathrm{TT}}
\newcommand{\afnode}[1]{\texttt{[#1]}}
\newcommand{\aftest}[1]{\texttt{tests/#1}}
\newcommand{\afdef}[1]{\texttt{def:#1}}
\newcommand{\afext}[1]{\texttt{ext:#1}}

\title{\bfseries Propagators of Linearised Fourth-Derivative Gravity:\\
\large Adversarial Proof Formalization and Verification Report}
\author{Generated by Claude Opus 4.6 via the AF (Adversarial Proof Framework)}
\date{5 March 2026}

\begin{document}
\maketitle

\begin{abstract}
We present the adversarial formalization and verification of the propagator
structure for linearised fourth-derivative gravity, with action
$I = \int d^4x\,[R^{(1)}_{\mu\nu}R^{(1)\mu\nu} - \beta(R^{(1)})^2]$.
The proof was formalized using the AF (Adversarial Proof Framework) CLI tool,
producing a proof tree of 38~nodes across 7~sections, supported by
10~formal definitions, 6~external references with ground-truth verification,
and 12~computer algebra test scripts (7~primary SymPy tests + 5~verifier-generated
scripts). Eight independent adversarial verifiers scrutinized every node;
9~challenges were raised and resolved, including one genuine sign convention
error (corrected) and one missing computational evidence file (created).
The final quality score is 100/100. All source materials, proof ledger entries,
and executable tests are included by reference.
\end{abstract}

\tableofcontents
\newpage

%% ===========================================================================
\section{Introduction}
\label{sec:intro}
%% ===========================================================================

This report documents the formal adversarial verification of the complete
propagator derivation for linearised fourth-derivative (quadratic curvature)
gravity. The main result is:

\begin{theorem}[Main Conjecture --- AF Node \afnode{1}]
\label{thm:main}
The propagators of linearised fourth-derivative gravity with action
\begin{equation}
I = \int d^4x\;\bigl[R^{(1)}_{\mu\nu}R^{(1)\mu\nu} - \beta\,(R^{(1)})^2\bigr]
\label{eq:action}
\end{equation}
decompose into scalar, vector, and tensor sectors with no cross-terms.
For $\beta \neq 1/3$, the scalar two-point functions are:
\begin{align}
\langle\Phi(k)\,\Phi(-k)\rangle &= \frac{3}{4k^4} + \frac{1}{2k^2p^2} + \frac{1-2\beta}{8(1-3\beta)\,p^4}, \label{eq:phiphi}\\
\langle\psi(k)\,\psi(-k)\rangle &= \frac{1-2\beta}{8(1-3\beta)\,p^4}, \label{eq:psipsi}\\
\langle\Phi(k)\,\psi(-k)\rangle &= -\frac{1}{4k^2p^2} - \frac{1-2\beta}{8(1-3\beta)\,p^4}; \label{eq:phipsi}
\end{align}
the vector propagator is
\begin{equation}
\langle V_i(k)\,V_j(-k)\rangle = \frac{P^T_{ij}(\mathbf{k})}{k^2\,p^2},\qquad
P^T_{ij} = \delta_{ij} - \frac{k_ik_j}{k^2}; \label{eq:vector}
\end{equation}
and the tensor propagator is
\begin{equation}
\langle h^{\TT}_{ij}(k)\,h^{\TT}_{kl}(-k)\rangle = \frac{2\,\Pi^{\TT}_{ijkl}(\mathbf{k})}{p^4}, \label{eq:tensor}
\end{equation}
where $\Pi^{\TT}_{ijkl} = \tfrac{1}{2}(P^T_{ik}P^T_{jl} + P^T_{il}P^T_{jk} - P^T_{ij}P^T_{kl})$.
\end{theorem}

Here $p^2 \equiv \omega^2 - \mathbf{k}^2$ is the Lorentzian four-momentum squared
(Definition~\afdef{lorentzian\_four\_momentum\_squared}, \S\ref{sec:defs}),
and we work in signature $(-,+,+,+)$.

\subsection{Methodology}

The proof was formalized using the \texttt{af} (Adversarial Proof Framework)
command-line tool, version \texttt{dev} (Go~1.25.5). The workflow comprises:

\begin{enumerate}
\item \textbf{Initialization}: \texttt{af init} creates a proof workspace with
  the conjecture as root node \afnode{1}.
\item \textbf{Refinement}: A prover agent decomposes the conjecture into 7~sections
  (nodes \afnode{1.1}--\afnode{1.7}), each further refined to leaf nodes.
\item \textbf{Definitions}: 10 formal definitions added via \texttt{af def-add}
  (Table~\ref{tab:defs}).
\item \textbf{External references}: 6 references added via \texttt{af add-external},
  with ground-truth string matches from source texts (Table~\ref{tab:externals}).
\item \textbf{Computational evidence}: 7 primary SymPy test scripts attached
  via \texttt{af attach} and validated via \texttt{af claim-test} (\S\ref{sec:tests}).
\item \textbf{Adversarial verification}: 8 independent verifier agents scrutinize
  every node, raising challenges for errors, gaps, or missing evidence (\S\ref{sec:verification}).
\end{enumerate}

The complete proof tree (38 nodes) is reproduced in Appendix~\ref{app:tree}.

%% ===========================================================================
\section{Proof Structure}
\label{sec:structure}
%% ===========================================================================

The proof tree has depth~3 and 38 nodes organized into 7 sections.
Each section corresponds to a logical stage of the derivation:

\begin{center}
\begin{tabular}{clcc}
\toprule
\textbf{Node} & \textbf{Section} & \textbf{Leaves} & \textbf{Tests} \\
\midrule
\afnode{1.1} & Action identification & 4 & --- \\
\afnode{1.2} & Gauge-invariant variables & 5 & --- \\
\afnode{1.3} & Scalar sector curvature & 4 & \aftest{test\_ricci\_scalar\_sector.py} \\
\afnode{1.4} & Vector \& tensor curvature & 5 & \aftest{test\_vector\_action.py}, \aftest{test\_tensor\_action.py} \\
\afnode{1.5} & Decomposed action & 4 & \aftest{test\_scalar\_action.py} \\
\afnode{1.6} & Determinant \& propagators & 4 & \aftest{test\_determinant.py}, \aftest{test\_matrix\_inverse.py} \\
\afnode{1.7} & Final results \& special cases & 4 & \aftest{test\_special\_cases.py} \\
\bottomrule
\end{tabular}
\end{center}

The logical dependencies are:
$\afnode{1.1} \to \afnode{1.2} \to \{\afnode{1.3}, \afnode{1.4}\} \to \afnode{1.5} \to \afnode{1.6} \to \afnode{1.7}$.

%% ===========================================================================
\section{Proof Content by Section}
\label{sec:content}
%% ===========================================================================

\subsection{Section 1: Action Identification (Nodes \afnode{1.1.1}--\afnode{1.1.4})}

The linearised Riemann tensor (\afdef{linearised\_riemann\_tensor}) is
\begin{equation}
R^{(1)}_{\mu\nu\rho\sigma} = \tfrac{1}{2}\bigl(
\p_\nu\p_\rho h_{\mu\sigma} + \p_\mu\p_\sigma h_{\nu\rho}
- \p_\mu\p_\rho h_{\nu\sigma} - \p_\nu\p_\sigma h_{\mu\rho}\bigr),
\label{eq:riemann}
\end{equation}
standard from Wald~\cite{Wald84} \S4.4 (\afext{Wald}; ground truth verified
via OCR from the \texttt{.djvu} source, pages~82--83, eq.~(4.4.3)--(4.4.4)).
Node \afnode{1.1.1} was verified by checking all four Riemann symmetries
($R_{abcd} = -R_{bacd}$, $R_{abcd} = -R_{abdc}$, $R_{abcd} = R_{cdab}$)
and the first Bianchi identity.

Contracting yields the linearised Ricci tensor (\afnode{1.1.2}, \afdef{linearised\_ricci\_tensor}):
\begin{equation}
2R^{(1)}_{\mu\nu} = \p_\lambda\p_\mu h^\lambda{}_\nu + \p_\lambda\p_\nu h^\lambda{}_\mu
- \p_\mu\p_\nu h - \dbox h_{\mu\nu},
\label{eq:ricci}
\end{equation}
cross-checked against Zee~\cite{Zee13} \S IX.4 eq.~(1), p.~563
(\afext{Zee}; ground truth: ``$R_{\mu\nu} = -\tfrac{1}{2}(\p^2 h_{\mu\nu}
- \p_\mu\p_\lambda h^\lambda{}_\nu - \p_\nu\p_\lambda h^\lambda{}_\mu
+ \p_\mu\p_\nu h^\lambda{}_\lambda) + O(h^2)$'').

The linearised Ricci scalar (\afnode{1.1.3}, \afdef{linearised\_ricci\_scalar}) is
$R^{(1)} = \p_\mu\p_\nu h^{\mu\nu} - \dbox h$.

\subsection{Section 2: Gauge-Invariant Variables (Nodes \afnode{1.2.1}--\afnode{1.2.5})}

The SVT decomposition (\afnode{1.2.1}, \afdef{svt\_decomposition},
\afext{SVT Decomposition (Bardeen 1980)}):
\begin{equation}
h_{00} = 2\phi,\quad
h_{0i} = \p_i B + S_i,\quad
h_{ij} = 2\psi\delta_{ij} + 2\p_i\p_j E + \p_i F_j + \p_j F_i + h^{\TT}_{ij},
\end{equation}
with $\p_i S_i = \p_i F_i = \p_i h^{\TT}_{ij} = h^{\TT}_{ii} = 0$.

The gauge transformation rule (\afnode{1.2.2}, amended during verification)
uses the \emph{Wald convention} $h_{\mu\nu} \to h_{\mu\nu} + \p_\mu\xi_\nu + \p_\nu\xi_\mu$
(eq.~(4.4.9) of \cite{Wald84}). This was the subject of the sole major
challenge during verification: the original formulation used a minus sign,
which was corrected (see \S\ref{sec:challenge122}).

The Bardeen potentials (\afnode{1.2.3}, \afdef{bardeen\_potentials}) are gauge-invariant:
$\Phi = \phi + \dot{B} - \ddot{E}$, $\Psi = \psi$, $V_i = S_i - \dot{F}_i$.
In the gauge $E = B = F_i = 0$, these reduce to the metric components (\afnode{1.2.4}).
The action~\eqref{eq:action}, built from linearised curvature, is gauge-invariant
(\afnode{1.2.5}, \afext{Gauge Invariance of Linearised Curvature}).

\subsection{Section 3: Scalar Sector Curvature (Nodes \afnode{1.3.1}--\afnode{1.3.4})}

Setting $h_{00} = 2\Phi$, $h_{0i} = 0$, $h_{ij} = 2\psi\delta_{ij}$
(scalar sector in the Bardeen gauge), we compute from~\eqref{eq:ricci}:
\begin{align}
R_{00} &= -\lap\Phi - 3\ddot\psi, \label{eq:R00}\\
R_{0i}\big|_S &= -2\,\p_i\dot\psi, \label{eq:R0i}\\
R_{ij}\big|_S &= \p_i\p_j(\Phi - \psi) - \dbox\psi\,\delta_{ij}, \label{eq:Rij}\\
R^{(1)} &= 2\lap\Phi - 4\lap\psi + 6\ddot\psi. \label{eq:Rscalar}
\end{align}
Each component was independently recalculated by Verifier-3 (\S\ref{sec:verification})
and confirmed by \aftest{test\_ricci\_scalar\_sector.py} (see \S\ref{sec:tests}).

\subsection{Section 4: Vector and Tensor Curvature (Nodes \afnode{1.4.1}--\afnode{1.4.5})}

\textbf{Vector sector} ($h_{0i} = V_i$, $\p_i V_i = 0$):
$R_{00} = 0$, $R_{0i} = -\tfrac{1}{2}\lap V_i$,
$R_{ij} = -\tfrac{1}{2}(\p_i\dot{V}_j + \p_j\dot{V}_i)$, $R = 0$.
Verified by \aftest{test\_vector\_action.py}.

\textbf{Tensor sector} ($h_{ij} = h^{\TT}_{ij}$):
$R_{ij} = -\tfrac{1}{2}\dbox h^{\TT}_{ij}$,
$R_{00} = R_{0i} = R = 0$ (by transversality and tracelessness).
Verified by \aftest{test\_tensor\_action.py}.

Critical check: the index raising $h^0{}_i = \eta^{00}h_{0i} = -V_i$ versus
$h^i{}_0 = \eta^{ij}h_{j0} = +V_i$ was verified to be handled correctly throughout
(Verifier-4).

\subsection{Section 5: Decomposed Action (Nodes \afnode{1.5.1}--\afnode{1.5.4})}

\begin{theorem}[Theorem 5.1 --- Node \afnode{1.5}]
The action decomposes as $I = I_{\TT} + I_V + I_S$ with no cross-terms, where:
\begin{align}
I_{\TT} &= \tfrac{1}{4}\int d^4x\;(\dbox\,h^{\TT}_{ij})^2, \\
I_V &= -\tfrac{1}{2}\int d^4x\; V_i\,\lap\dbox\, V_i, \\
I_S &= \int \frac{d^4k}{(2\pi)^4}\;
\begin{pmatrix}\Phi^* & \psi^*\end{pmatrix} M(k)
\begin{pmatrix}\Phi \\ \psi\end{pmatrix},
\end{align}
with the scalar kinetic matrix (\afdef{scalar\_action\_matrix})
\begin{equation}
M = \begin{pmatrix}
2(1-2\beta)\,k^4 & 4(1-3\beta)\,k^2p^2 + 2(1-2\beta)\,k^4 \\
4(1-3\beta)\,k^2p^2 + 2(1-2\beta)\,k^4 &
12(1-3\beta)\,p^4 + 8(1-3\beta)\,k^2p^2 + 2(1-2\beta)\,k^4
\end{pmatrix}.
\label{eq:matrix_M}
\end{equation}
\end{theorem}

The absence of cross-terms (\afnode{1.5.1}) follows from Schur orthogonality
for $\mathrm{SO}(3)$ irreducible representations
(\afext{Schur Orthogonality for SO(3) Irreps}):
the scalar, vector, and tensor sectors transform as spin-0, spin-1, and spin-2.
Verifier-5 confirmed this argument is rigorous even though the operator involves
time derivatives, since $\p_t$ commutes with spatial rotations and acts as a scalar
under $\mathrm{SO}(3)$.

The matrix $M$ was derived by substituting~\eqref{eq:R00}--\eqref{eq:Rscalar}
into $R_{\mu\nu}R^{\mu\nu} - \beta R^2$ and passing to Fourier space
($\lap \to -k^2$, $\dbox \to p^2$). This derivation is verified by
\aftest{test\_scalar\_action.py}, which independently recomputes all entries
of $M$ from the Ricci components.

\subsection{Section 6: Determinant and Propagators (Nodes \afnode{1.6.1}--\afnode{1.6.4})}

\begin{lemma}[Lemma 6.1 --- Node \afnode{1.6.2}]
$\det M = 8(1-3\beta)\,k^4\,p^4$.
\end{lemma}

The key algebraic identity is $3(1-2\beta) - 2(1-3\beta) = 1$, which causes the
$k^8$ and $k^6 p^2$ terms to cancel in $M_{11}M_{22} - M_{12}^2$.
Verified by \aftest{test\_determinant.py}.

The propagator normalization (\afnode{1.6.3}): for the action
$I_S = \int \Phi^*_a M_{ab} \Phi_b$ integrated over all $k$, real fields satisfy
$\Phi_a(-k) = \Phi^*_a(k)$, giving canonical form $\tfrac{1}{2}\int \Phi(-k)\,A(k)\,\Phi(k)$
with $A = 2M$. Hence $G = A^{-1} = \tfrac{1}{2}M^{-1}$.
This is consistent across all sectors:
\begin{itemize}
\item Tensor: $I_{\TT} = \tfrac{1}{4}\int p^4|h^{\TT}|^2$ gives $G = 2/p^4$ \checkmark
\item Vector: $I_V = \tfrac{1}{2}\int k^2p^2|V|^2$ gives $G = 1/(k^2p^2)$ \checkmark
\end{itemize}

The scalar propagator expressions~\eqref{eq:phiphi}--\eqref{eq:phipsi}
follow from $\tfrac{1}{2}M^{-1}$ by partial fractions (\afnode{1.6.4}).
Verified by \aftest{test\_matrix\_inverse.py}, which also confirms
$2M \cdot \tfrac{1}{2}M^{-1} = I$.

\subsection{Section 7: Vector/Tensor Propagators and Special Cases (Nodes \afnode{1.7.1}--\afnode{1.7.4})}

The vector propagator~\eqref{eq:vector} restricts $1/(k^2p^2)$ to the transverse
subspace via $P^T_{ij}$ (\afnode{1.7.1}, \afdef{transverse\_projector}).
The tensor propagator~\eqref{eq:tensor} uses the TT projector
(\afnode{1.7.2}, \afdef{tt\_projector}).
Verifier-7 confirmed $\Pi^{\TT}$ is idempotent, traceless, and transverse,
with $\mathrm{Tr}(\Pi^{\TT}) = 2$ (two polarizations).

\begin{remark}[Conformal gravity --- Node \afnode{1.7.4}]
At $\beta = 1/3$, $\det M = 0$. The action reduces to
$\tfrac{1}{2}C_{\mu\nu\rho\sigma}C^{\mu\nu\rho\sigma}$ (up to the Gauss--Bonnet
total derivative), and conformal symmetry removes one scalar DOF.
Verified by \aftest{test\_special\_cases.py} and the independent script
\aftest{verify\_weyl\_identity\_independent.py}. The identity
$R_{\mu\nu}R^{\mu\nu} - \tfrac{1}{3}R^2 = \tfrac{1}{2}C^2 - \tfrac{1}{2}\mathcal{G}$
(where $\mathcal{G}$ is the Gauss--Bonnet term) was confirmed using the 4D
Gauss--Bonnet relation (\afext{Gauss-Bonnet at Linear Order}).
\end{remark}

\begin{remark}[Node \afnode{1.7.4}]
At $\beta = 1/2$: $(1-2\beta) = 0$, so $\langle\psi\psi\rangle = 0$
and all $p^{-4}$ poles vanish from the scalar propagators.
\end{remark}

%% ===========================================================================
\section{Definitions}
\label{sec:defs}
%% ===========================================================================

\begin{longtable}{lp{10cm}}
\toprule
\textbf{Name} & \textbf{Content} \\
\midrule
\endhead
\texttt{lorentzian\_four\_momentum\_squared} &
$p^2 \equiv \omega^2 - \mathbf{k}^2$; $\dbox \leftrightarrow p^2$ in Fourier space \\
\texttt{fourth\_derivative\_gravity\_action} &
$I = \int d^4x\,[R^{(1)}_{\mu\nu}R^{(1)\mu\nu} - \beta(R^{(1)})^2]$ \\
\texttt{linearised\_riemann\_tensor} &
$R^{(1)}_{\mu\nu\rho\sigma} = \tfrac{1}{2}(\p_\nu\p_\rho h_{\mu\sigma} + \cdots)$ \\
\texttt{linearised\_ricci\_tensor} &
$2R^{(1)}_{\mu\nu} = \p_\lambda\p_\mu h^\lambda{}_\nu + \cdots$ \\
\texttt{linearised\_ricci\_scalar} &
$R^{(1)} = \p_\mu\p_\nu h^{\mu\nu} - \dbox h$ \\
\texttt{svt\_decomposition} &
$h_{00} = 2\phi$, $h_{0i} = \p_i B + S_i$, $h_{ij} = 2\psi\delta_{ij} + \cdots$ \\
\texttt{bardeen\_potentials} &
$\Phi = \phi + \dot{B} - \ddot{E}$, $\Psi = \psi$, $V_i = S_i - \dot{F}_i$ \\
\texttt{transverse\_projector} &
$P^T_{ij} = \delta_{ij} - k_ik_j/k^2$ \\
\texttt{tt\_projector} &
$\Pi^{\TT}_{ijkl} = \tfrac{1}{2}(P^T_{ik}P^T_{jl} + P^T_{il}P^T_{jk} - P^T_{ij}P^T_{kl})$ \\
\texttt{scalar\_action\_matrix} &
The $2\times 2$ matrix $M$ in eq.~\eqref{eq:matrix_M} \\
\bottomrule
\caption{Formal definitions registered in the AF proof via \texttt{af def-add}.}
\label{tab:defs}
\end{longtable}

%% ===========================================================================
\section{External References}
\label{sec:externals}
%% ===========================================================================

\begin{longtable}{lp{8.5cm}p{3cm}}
\toprule
\textbf{Name} & \textbf{Source} & \textbf{Ground Truth} \\
\midrule
\endhead
Wald \S4.4 &
Wald, R.M.\ (1984). \textit{General Relativity}. U.\ Chicago Press, \S4.4 &
OCR from \texttt{.djvu}, pp.~82--83, eq.~(4.4.3)--(4.4.4) \\
Zee IX.4 &
Zee, A.\ (2013). \textit{Einstein Gravity in a Nutshell}. Princeton, \S IX.4, p.~563, eq.~(1) &
pdftotext from \texttt{.pdf}: ``$R_{\mu\nu} = -\tfrac{1}{2}(\p^2 h_{\mu\nu} - \cdots)$'' \\
Bardeen 1980 &
Bardeen, J.M.\ (1980). Phys.\ Rev.\ D, 22(8):1882 &
Standard SVT decomposition \\
Gauge Invariance &
Standard: $\delta R^{(1)}_{\mu\nu\rho\sigma} = 0$ under linearised diffeos &
Textbook result \\
Schur Orthogonality &
Sternberg, S.\ (1994). \textit{Group Theory and Physics}. Cambridge &
$\mathrm{SO}(3)$ irrep orthogonality \\
Gauss--Bonnet &
$R_{\mu\nu}^2 - \tfrac{1}{3}R^2 = \tfrac{1}{2}C^2$ mod total derivs in 4D &
Standard higher-deriv.\ gravity \\
\bottomrule
\caption{External references registered via \texttt{af add-external}, with ground-truth verification method.}
\label{tab:externals}
\end{longtable}

%% ===========================================================================
\section{Computer Algebra Tests}
\label{sec:tests}
%% ===========================================================================

All tests are executable Python scripts using SymPy~1.14.0 (installed in
\texttt{/tmp/cas-venv/}). Each script exits with code~0 on success.
Tests were attached to proof nodes via \texttt{af attach} and validated
via \texttt{af claim-test}.

\subsection{Primary Tests (7 scripts)}

\begin{longtable}{llp{6.5cm}l}
\toprule
\textbf{File} & \textbf{Node} & \textbf{Verifies} & \textbf{LOC} \\
\midrule
\endhead
\texttt{test\_determinant.py} & \afnode{1.6.2} &
$\det M = 8(1-3\beta)k^4p^4$; spot checks at $\beta = 0, 1/3, 1/2$ & 74 \\
\texttt{test\_matrix\_inverse.py} & \afnode{1.6.4} &
All 3 scalar propagators via $\tfrac{1}{2}M^{-1}$; $2M \cdot G = I$ & 128 \\
\texttt{test\_ricci\_scalar\_sector.py} & \afnode{1.3} &
$R_{00}$, $R_{0i}$, $R_{ij}$, $R$ in scalar gauge; dual method for $R$ & 167 \\
\texttt{test\_scalar\_action.py} & \afnode{1.5.4} &
Full derivation of $M$ from $R_{\mu\nu}R^{\mu\nu} - \beta R^2$ in Fourier & 211 \\
\texttt{test\_vector\_action.py} & \afnode{1.5.3} &
$I_V = -\tfrac{1}{2}\int V_i\lap\dbox V_i$; Fourier: $\tfrac{1}{2}k^2p^2|V|^2$ & 198 \\
\texttt{test\_tensor\_action.py} & \afnode{1.5.2} &
$I_{\TT} = \tfrac{1}{4}\int(\dbox h^{\TT}_{ij})^2$; $\beta$-independence & 194 \\
\texttt{test\_special\_cases.py} & \afnode{1.7.4} &
$\beta = 1/2$ (reduced DOF), $\beta = 1/3$ (conformal), sum rule & 181 \\
\bottomrule
\caption{Primary SymPy verification scripts, attached as computational evidence.}
\label{tab:tests}
\end{longtable}

\subsection{Verifier-Generated Scripts (5 scripts)}

During adversarial verification, Verifier-7 independently wrote additional
tests stored in \texttt{tests/}:

\begin{itemize}
\item \texttt{verify\_1\_7\_1.py} --- $P^T_{ij}$ idempotency, transversality, trace
\item \texttt{verify\_1\_7\_2.py} --- $\Pi^{\TT}_{ijkl}$ idempotency, tracelessness,
  consistency with polarization modes
\item \texttt{verify\_1\_7\_3.py} --- $1/p^4$ pole structure, instantaneous $1/k^4$ piece
\item \texttt{verify\_1\_7\_4.py} --- $\beta = 1/2$, $\beta = 1/3$ limits, rank of $M$
\item \texttt{verify\_weyl\_identity\_independent.py} --- Algebraic verification of
  $R_{\mu\nu}^2 - \tfrac{1}{3}R^2 = \tfrac{1}{2}C^2 - \tfrac{1}{2}\mathcal{G}$
\end{itemize}

%% ===========================================================================
\section{Adversarial Verification}
\label{sec:verification}
%% ===========================================================================

Verification proceeded in two batches of 4~verifiers each, processing the proof
bottom-up (leaf nodes first, parents after children).

\subsection{Batch 1: Sections 1--4}

\begin{longtable}{llccp{5.5cm}}
\toprule
\textbf{Agent} & \textbf{Section} & \textbf{Nodes} & \textbf{Challenges} & \textbf{Key Checks} \\
\midrule
\endhead
Verifier-1 & \S1 Action & 5 & 0 &
All 4 Riemann symmetries + Bianchi identity; derivation from Christoffel symbols \\
Verifier-2 & \S2 Gauge & 6 & \textbf{4} &
\textbf{Major}: sign convention error in \afnode{1.2.2} (see \S\ref{sec:challenge122}).
3 minor/notes on dependencies and notation. \\
Verifier-3 & \S3 Scalar & 5 & 0 &
Full independent recalculation of all 4 Ricci components + SymPy cross-check \\
Verifier-4 & \S4 Vec+Tens & 6 & 0 &
Index raising $h^0{}_i = -V_i$ vs.\ $h^i{}_0 = +V_i$; SymPy confirmation \\
\bottomrule
\caption{Batch 1 verification results.}
\end{longtable}

\subsection{Batch 2: Sections 5--7 and Root}

\begin{longtable}{llccp{5.5cm}}
\toprule
\textbf{Agent} & \textbf{Section} & \textbf{Nodes} & \textbf{Challenges} & \textbf{Key Checks} \\
\midrule
\endhead
Verifier-5 & \S5 Action & 5 & \textbf{3} &
\textbf{Major}: missing \texttt{test\_scalar\_action.py} (created, see \S\ref{sec:challenge154}).
2 notes on $\dbox$ sign convention in Definition~0. \\
Verifier-6 & \S6 Det+Prop & 5 & 0 &
Normalization $G = \tfrac{1}{2}M^{-1}$; reality condition $M(-k) = M(k)$;
partial fractions by hand for $\beta = 0$ \\
Verifier-7 & \S7 Results & 5 & 1 &
$\Pi^{\TT}$ idempotency; Weyl$^2$ identity via Gauss--Bonnet;
1 note on Definition~0 sign \\
Verifier-8 & Root & 1+6 & 1 &
All 6 tests pass; logical flow verified across sections;
normalization consistency; 1 minor note \\
\bottomrule
\caption{Batch 2 verification results.}
\end{longtable}

\subsection{Challenge: Sign Convention Error (Node \afnode{1.2.2})}
\label{sec:challenge122}

\textbf{Challenge ID}: \texttt{ch-54ce4583716368f0} (severity: major, blocking).

\textbf{Issue}: Node \afnode{1.2.2} originally stated the gauge transformation as
$h_{\mu\nu} \to h_{\mu\nu} - \p_\mu\xi_\nu - \p_\nu\xi_\mu$, but derived
transformation rules ($\phi \to \phi - \dot{T}$, etc.)\ consistent with the
\emph{opposite} sign convention $h_{\mu\nu} \to h_{\mu\nu} + \p_\mu\xi_\nu + \p_\nu\xi_\mu$
(Wald eq.~(4.4.9)).

\textbf{Analysis}: With $\xi_0 = \eta_{00}\xi^0 = -T$ in signature $(-,+,+,+)$,
the minus-sign formula gives $\delta h_{00} = -2\p_0\xi_0 = +2\dot{T}$, hence
$\phi \to \phi + \dot{T}$---contradicting the stated (and correct)
$\phi \to \phi - \dot{T}$.

\textbf{Resolution}: The node statement was amended via \texttt{af amend} to use
the Wald convention with plus sign, with an explicit derivation showing
$\p_0\xi_0 = -\dot{T}$, $h_{00} \to h_{00} + 2\p_0\xi_0 = 2\phi - 2\dot{T}$.

\textbf{Impact}: The error was confined to the explanatory formula; all downstream
transformation \emph{rules} and \emph{results} were already correct.

\subsection{Challenge: Missing Test File (Node \afnode{1.5.4})}
\label{sec:challenge154}

\textbf{Challenge ID}: \texttt{ch-84d79921ac409f9a} (severity: major, blocking).

\textbf{Issue}: Node \afnode{1.5.4} claimed ``Verified by SymPy
\texttt{test\_scalar\_action.py}'' but the file did not exist.

\textbf{Resolution}: The test script was created, deriving the matrix $M$ from
scratch via Fourier-space computation of $R_{\mu\nu}R^{\mu\nu} - \beta R^2$.
All entries verified. The script was attached via \texttt{af attach} and
validated via \texttt{af claim-test}. The test includes cross-checks at
$\beta = 0, 1/3, 1/2$ and on-shell ($p^2 = 0$).

\subsection{Minor Issues}

A notation inconsistency in Definition~0 was identified by three verifiers
independently: the definition states $\dbox \leftrightarrow -p^2$ with
$p^2 = \omega^2 - k^2$, but the correct Fourier relation is
$\dbox \leftrightarrow +p^2$. This does not affect any computations (all
calculations use the correct signs).

%% ===========================================================================
\section{Final Metrics}
\label{sec:metrics}
%% ===========================================================================

\begin{center}
\begin{tabular}{lr}
\toprule
\textbf{Metric} & \textbf{Value} \\
\midrule
Total proof nodes & 38 \\
Validated & 38 (100\%) \\
Pending/Admitted/Refuted & 0 / 0 / 0 \\
Max tree depth & 3 \\
Definitions & 10 \\
External references & 6 \\
Challenges raised & 9 \\
Challenges resolved & 9 \\
Challenge density & 0.24 per node \\
Primary test scripts & 7 \\
Verifier-generated scripts & 5 \\
Total test LOC & $\sim$1150 \\
Quality score & \textbf{100/100 (EXCELLENT)} \\
\bottomrule
\end{tabular}
\end{center}

%% ===========================================================================
\section{Conclusion}
\label{sec:conclusion}
%% ===========================================================================

The propagator structure of linearised fourth-derivative gravity has been
formalized as a 38-node adversarial proof tree, verified by 8~independent
agents, and supported by 12~computer algebra scripts. One genuine mathematical
error (sign convention) and one evidentiary gap (missing test) were discovered
and corrected during verification. The final proof is validated at 100\%
with quality score 100/100.

All materials are available in the repository at
\texttt{examples12/}:
\begin{itemize}
\item Proof ledger: \texttt{ledger/000001.json}--\texttt{000245.json} (245 entries)
\item Definitions: \texttt{defs/} (10 JSON files)
\item External references: \texttt{externals/} (6 JSON files)
\item Test scripts: \texttt{tests/} (12 Python files)
\item Ground-truth source: \texttt{refs/Zee2013.pdf} (local copy)
\item Exported proof tree: \texttt{proof\_tree.tex} (Appendix~\ref{app:tree})
\item Original derivation: \texttt{fourth\_derivative\_gravity\_propagators.md}
\end{itemize}

\begin{thebibliography}{9}
\bibitem{Wald84} R.M.~Wald, \textit{General Relativity},
  University of Chicago Press, 1984.
\bibitem{Zee13} A.~Zee, \textit{Einstein Gravity in a Nutshell},
  Princeton University Press, 2013.
\bibitem{Bardeen80} J.M.~Bardeen, ``Gauge-invariant cosmological perturbations,''
  Phys.\ Rev.\ D \textbf{22}, 1882 (1980).
\bibitem{Stelle77} K.S.~Stelle, ``Renormalization of higher-derivative quantum gravity,''
  Phys.\ Rev.\ D \textbf{16}, 953 (1977).
\end{thebibliography}

%% ===========================================================================
\appendix
\section{Exported Proof Tree}
\label{app:tree}
%% ===========================================================================

The following is the complete AF proof tree exported via
\texttt{af export --format latex}. All 38~nodes are shown with their
statements, types, inference rules, epistemic states, and taint states.
Node identifiers (e.g., \afnode{1.3.2}) correspond to references throughout
this report.

\bigskip
\hrule
\bigskip


\section*{Node 1}

\textbf{Statement:} The propagators of linearised fourth-derivative gravity (action I = ∫d⁴x [R\textasciicircum{}(1)\_\{μν\}R\textasciicircum{}\{(1)μν\} - β(R\textasciicircum{}(1))²]) decompose into scalar, vector, and tensor sectors with no cross-terms. For β ≠ 1/3, the scalar propagators are ⟨ΦΦ⟩ = 3/(4k⁴) + 1/(2k²p²) + (1-2β)/[8(1-3β)p⁴], ⟨ψψ⟩ = (1-2β)/[8(1-3β)p⁴], ⟨Φψ⟩ = -1/(4k²p²) - (1-2β)/[8(1-3β)p⁴]; the vector propagator is ⟨V\_iV\_j⟩ = P\textasciicircum{}T\_\{ij\}/(k²p²); and the tensor propagator is ⟨h\textasciicircum{}\{TT\}\_\{ij\}h\textasciicircum{}\{TT\}\_\{kl\}⟩ = 2Π\textasciicircum{}\{TT\}\_\{ijkl\}/p⁴.

\textbf{Type:} claim

\textbf{Inference:} assumption

\textbf{Status:} validated

\textbf{Taint:} clean

\begin{enumerate}
\item \textbf{Node 1.1:} Claim 1.1: The integrand equals R\textasciicircum{}(1)\_\{μν\}R\textasciicircum{}\{(1)μν\} - β(R\textasciicircum{}(1))², where R\textasciicircum{}(1)\_\{μν\} is the linearised Ricci tensor and R\textasciicircum{}(1)\} the linearised Ricci scalar. The linearised Riemann tensor is R\textasciicircum{}(1)\_\{μνρσ\} = (1/2)(∂\_ν∂\_ρ h\_\{μσ\} + ∂\_μ∂\_σ h\_\{νρ\} - ∂\_μ∂\_ρ h\_\{νσ\} - ∂\_ν∂\_σ h\_\{μρ\}). Contracting gives 2R\textasciicircum{}(1)\_\{μν\} = ∂\_λ∂\_μ h\textasciicircum{}λ\_ν + ∂\_λ∂\_ν h\textasciicircum{}λ\_μ - ∂\_μ∂\_ν h - □h\_\{μν\}, and R\textasciicircum{}(1) = ∂\_μ∂\_ν h\textasciicircum{}\{μν\} - □h. [Refs: Wald §4.4, Zee IX.4 eq.(1)]

  Type: claim, Inference: assumption, Status: validated, Taint: clean

\begin{enumerate}
\item \textbf{Node 1.1.1:} 1.1.1: The linearised Riemann tensor for g\_\{μν\} = η\_\{μν\} + h\_\{μν\} is R\textasciicircum{}(1)\_\{μνρσ\} = (1/2)(∂\_ν∂\_ρ h\_\{μσ\} + ∂\_μ∂\_σ h\_\{νρ\} - ∂\_μ∂\_ρ h\_\{νσ\} - ∂\_ν∂\_σ h\_\{μρ\}). [Standard, Wald §4.4]

  Type: claim, Inference: assumption, Status: validated, Taint: clean

\item \textbf{Node 1.1.2:} 1.1.2: Contracting R\textasciicircum{}(1)\_\{μρν\}\textasciicircum{}ρ gives 2R\textasciicircum{}(1)\_\{μν\} = ∂\_λ∂\_μ h\textasciicircum{}λ\_ν + ∂\_λ∂\_ν h\textasciicircum{}λ\_μ - ∂\_μ∂\_ν h - □h\_\{μν\} where h = η\textasciicircum{}\{μν\}h\_\{μν\}. [Contract 1.1.1 with η\textasciicircum{}\{νσ\}, relabel]

  Type: claim, Inference: assumption, Status: validated, Taint: clean

\item \textbf{Node 1.1.3:} 1.1.3: The linearised Ricci scalar is R\textasciicircum{}(1) = η\textasciicircum{}\{μν\}R\textasciicircum{}(1)\_\{μν\} = ∂\_μ∂\_ν h\textasciicircum{}\{μν\} - □h. [Trace of 1.1.2]

  Type: claim, Inference: assumption, Status: validated, Taint: clean

\item \textbf{Node 1.1.4:} 1.1.4: The quadratic invariant R\textasciicircum{}(1)\_\{μν\}R\textasciicircum{}\{(1)μν\} is formed by contracting (1/4)(2R\textasciicircum{}(1)\_\{μν\})(2R\textasciicircum{}\{(1)μν\}), yielding R\textasciicircum{}(1)\_\{μν\}R\textasciicircum{}\{(1)μν\}. The β-term is β(R\textasciicircum{}(1))². Their difference gives the stated integrand. QED.

  Type: claim, Inference: assumption, Status: validated, Taint: clean

\end{enumerate}
\item \textbf{Node 1.2:} Claim 2.1: In the gauge E=0, B=0, F\_i=0, the residual metric components identify with Bardeen variables: φ=Φ, S\_i=V\_i, ψ=ψ, h\textasciicircum{}\{TT\}\_\{ij\}=h\textasciicircum{}\{TT\}\_\{ij\}. Since I is built from linearised curvature (gauge-invariant at linear order), the result in Bardeen variables holds in any gauge. [Refs: SVT Decomposition (Bardeen 1980), Gauge Invariance of Linearised Curvature]

  Type: claim, Inference: assumption, Status: validated, Taint: clean

\begin{enumerate}
\item \textbf{Node 1.2.1:} 2.1.1: The SVT decomposition of the metric perturbation is h\_\{00\}=2φ, h\_\{0i\}=∂\_iB+S\_i, h\_\{ij\}=2ψδ\_\{ij\}+2∂\_i∂\_jE+∂\_iF\_j+∂\_jF\_i+h\textasciicircum{}\{TT\}\_\{ij\} with ∂\_iS\_i=0, ∂\_iF\_i=0, ∂\_ih\textasciicircum{}\{TT\}\_\{ij\}=0, h\textasciicircum{}\{TT\}\_\{ii\}=0. [Standard SVT decomposition, Ref: SVT Decomposition (Bardeen 1980)]

  Type: claim, Inference: assumption, Status: validated, Taint: clean

\item \textbf{Node 1.2.2:} 2.1.2: Under x\textasciicircum{}μ→x\textasciicircum{}μ+ξ\textasciicircum{}μ with ξ\textasciicircum{}0=T, ξ\textasciicircum{}i=∂\textasciicircum{}iL+L\textasciicircum{}i\_T (∂\_iL\textasciicircum{}i\_T=0), the gauge transformation h\_\{μν\}→h\_\{μν\}+∂\_μξ\_ν+∂\_νξ\_μ (Wald convention, eq.4.4.9) gives: φ→φ-Ṫ, B→B+T-L̇, E→E-L, ψ→ψ, S\_i→S\_i-L̇\textasciicircum{}T\_i, F\_i→F\_i-L\textasciicircum{}T\_i. [Here ξ\_0=η\_\{00\}ξ\textasciicircum{}0=-T with (-,+,+,+), so ∂\_0ξ\_0=-Ṫ and h\_\{00\}→h\_\{00\}+2∂\_0ξ\_0=2φ-2Ṫ, giving φ→φ-Ṫ.]

  Type: claim, Inference: assumption, Status: validated, Taint: clean

\item \textbf{Node 1.2.3:} 2.1.3: The Bardeen potentials Φ=φ+Ḃ-Ë and Ψ=ψ are gauge-invariant. Verification: Ḃ-Ë→Ḃ+Ṫ-L̈-Ë+L̈=Ḃ-Ë+Ṫ, so Φ→φ-Ṫ+Ḃ-Ë+Ṫ=Φ. Vector: V\_i=S\_i-Ḟ\_i→S\_i-L̇\textasciicircum{}T\_i-Ḟ\_i+L̇\textasciicircum{}T\_i=V\_i. ✓

  Type: claim, Inference: assumption, Status: validated, Taint: clean

\item \textbf{Node 1.2.4:} 2.1.4: In the gauge E=B=F\_i=0: Φ=φ, V\_i=S\_i, Ψ=ψ, h\textasciicircum{}\{TT\}\_\{ij\}=h\textasciicircum{}\{TT\}\_\{ij\}. [Direct substitution into 2.1.3]

  Type: claim, Inference: assumption, Status: validated, Taint: clean

\item \textbf{Node 1.2.5:} 2.1.5: The action I is constructed from R\textasciicircum{}(1)\_\{μν\} and R\textasciicircum{}(1)\}, both gauge-invariant at linear order (the linearised Riemann tensor is invariant under h\_\{μν\}→h\_\{μν\}+∂\_μξ\_ν+∂\_νξ\_μ). Therefore I expressed in Bardeen variables is gauge-independent. QED. [Ref: Gauge Invariance of Linearised Curvature]

  Type: claim, Inference: assumption, Status: validated, Taint: clean

\end{enumerate}
\item \textbf{Node 1.3:} Claim 3.1: In the scalar sector gauge (h\_\{00\}=2Φ, h\_\{0i\}=0, h\_\{ij\}=2ψδ\_\{ij\}): R\_\{00\} = -∇²Φ - 3ψ̈, R\_\{0i\}|\_S = -2∂\_iψ̇, R\_\{ij\}|\_S = ∂\_i∂\_j(Φ-ψ) - □ψ δ\_\{ij\}, R\textasciicircum{}(1) = 2∇²Φ - 4∇²ψ + 6ψ̈. [Computed from linearised\_ricci\_tensor definition]

  Type: claim, Inference: assumption, Status: validated, Taint: clean

\begin{enumerate}
\item \textbf{Node 1.3.1:} 3.1.1 (R\_\{00\}): Set h\_\{00\}=2Φ, h\_\{0i\}=0, h\_\{ij\}=2ψδ\_\{ij\}. Then h=-2Φ+6ψ. From 1.1.2: 2R\_\{00\}=2∂\_λ∂\_0h\textasciicircum{}λ\_0-∂\_0²h-□h\_\{00\}. With h\textasciicircum{}0\_0=-2Φ, h\textasciicircum{}i\_0=0: ∂\_λ∂\_0h\textasciicircum{}λ\_0=-2Φ̈. Also ∂\_0²h=-2Φ̈+6ψ̈, □h\_\{00\}=2□Φ. Assembling: 2R\_\{00\}=-4Φ̈-(-2Φ̈+6ψ̈)-2□Φ=-2Φ̈-6ψ̈-2□Φ. Since □Φ=-Φ̈+∇²Φ: 2R\_\{00\}=-2∇²Φ-6ψ̈. Hence R\_\{00\}=-∇²Φ-3ψ̈. ✓

  Type: claim, Inference: assumption, Status: validated, Taint: clean

\item \textbf{Node 1.3.2:} 3.1.2 (R\_\{0i\}): 2R\_\{0i\}=∂\_λ∂\_0h\textasciicircum{}λ\_i+∂\_λ∂\_ih\textasciicircum{}λ\_0-∂\_0∂\_ih-□h\_\{0i\}. With h\textasciicircum{}j\_i=2ψδ\_\{ji\}, h\textasciicircum{}0\_i=0: ∂\_λ∂\_0h\textasciicircum{}λ\_i=2∂\_iψ̇. ∂\_λ∂\_ih\textasciicircum{}λ\_0=-2∂\_iΦ̇. ∂\_0∂\_ih=-2∂\_iΦ̇+6∂\_iψ̇, □h\_\{0i\}=0. Thus 2R\_\{0i\}=2∂\_iψ̇-2∂\_iΦ̇-(-2∂\_iΦ̇+6∂\_iψ̇)=-4∂\_iψ̇. Hence R\_\{0i\}=-2∂\_iψ̇. ✓

  Type: claim, Inference: assumption, Status: validated, Taint: clean

\item \textbf{Node 1.3.3:} 3.1.3 (R\_\{ij\}): 2R\_\{ij\}=∂\_λ∂\_ih\textasciicircum{}λ\_j+∂\_λ∂\_jh\textasciicircum{}λ\_i-∂\_i∂\_jh-□h\_\{ij\}. ∂\_λ∂\_ih\textasciicircum{}λ\_j=∂\_k∂\_i(2ψδ\_\{kj\})=2∂\_i∂\_jψ, symmetrically for second term: total 4∂\_i∂\_jψ. ∂\_i∂\_jh=∂\_i∂\_j(-2Φ+6ψ), □h\_\{ij\}=2□ψδ\_\{ij\}. Combining: 2R\_\{ij\}=4∂\_i∂\_jψ+2∂\_i∂\_jΦ-6∂\_i∂\_jψ-2□ψδ\_\{ij\}=2∂\_i∂\_j(Φ-ψ)-2□ψδ\_\{ij\}. Hence R\_\{ij\}=∂\_i∂\_j(Φ-ψ)-□ψδ\_\{ij\}. ✓

  Type: claim, Inference: assumption, Status: validated, Taint: clean

\item \textbf{Node 1.3.4:} 3.1.4 (R\textasciicircum{}(1)): Using R\textasciicircum{}(1)=∂\_μ∂\_νh\textasciicircum{}\{μν\}-□h. ∂\_μ∂\_νh\textasciicircum{}\{μν\}=∂\_0²h\textasciicircum{}\{00\}+∂\_i∂\_jh\textasciicircum{}\{ij\}=2Φ̈+2∇²ψ. □h=(-∂\_t²+∇²)(-2Φ+6ψ)=2Φ̈-6ψ̈-2∇²Φ+6∇²ψ. Therefore R\textasciicircum{}(1)=2Φ̈+2∇²ψ-2Φ̈+6ψ̈+2∇²Φ-6∇²ψ=2∇²Φ-4∇²ψ+6ψ̈. ✓ QED.

  Type: claim, Inference: assumption, Status: validated, Taint: clean

\end{enumerate}
\item \textbf{Node 1.4:} Claims 4.1-4.2: Vector sector (h\_\{0i\}=V\_i, ∂\_iV\_i=0): R\_\{00\}=0, R\_\{0i\}=-(1/2)∇²V\_i, R\_\{ij\}=-(1/2)(∂\_iV̇\_j+∂\_jV̇\_i), R=0. Tensor sector (h\_\{ij\}=h\textasciicircum{}\{TT\}\_\{ij\}): R\_\{ij\}=-(1/2)□h\textasciicircum{}\{TT\}\_\{ij\}, R\_\{00\}=R\_\{0i\}=R=0. [Computed from linearised\_ricci\_tensor with transversality and tracelessness]

  Type: claim, Inference: assumption, Status: validated, Taint: clean

\begin{enumerate}
\item \textbf{Node 1.4.1:} 4.1.1 (Vector R\_\{00\}): Set h\_\{00\}=0, h\_\{0i\}=V\_i, h\_\{ij\}=0, h=0. 2R\_\{00\}=2∂\_λ∂\_0h\textasciicircum{}λ\_0-□h\_\{00\}. h\textasciicircum{}0\_0=0, h\textasciicircum{}i\_0=V\_i. ∂\_λ∂\_0h\textasciicircum{}λ\_0=∂\_i∂\_0V\_i=∂\_0(∂\_iV\_i)=0 (transversality). □h\_\{00\}=0. Hence R\_\{00\}=0. ✓

  Type: claim, Inference: assumption, Status: validated, Taint: clean

\item \textbf{Node 1.4.2:} 4.1.2 (Vector R\_\{0i\}): 2R\_\{0i\}=∂\_λ∂\_0h\textasciicircum{}λ\_i+∂\_λ∂\_ih\textasciicircum{}λ\_0-□h\_\{0i\}. h\textasciicircum{}0\_i=-V\_i, h\textasciicircum{}j\_i=0. First term: ∂\_0²(-V\_i)=-V̈\_i. Second: ∂\_j∂\_iV\_j=0 (transversality). 2R\_\{0i\}=-V̈\_i-□V\_i=-V̈\_i+V̈\_i-∇²V\_i=-∇²V\_i. Hence R\_\{0i\}=-(1/2)∇²V\_i. ✓

  Type: claim, Inference: assumption, Status: validated, Taint: clean

\item \textbf{Node 1.4.3:} 4.1.3 (Vector R\_\{ij\}): 2R\_\{ij\}=∂\_λ∂\_ih\textasciicircum{}λ\_j+∂\_λ∂\_jh\textasciicircum{}λ\_i-∂\_i∂\_jh-□h\_\{ij\}. h\textasciicircum{}0\_j=-V\_j, h\textasciicircum{}k\_j=0. ∂\_λ∂\_ih\textasciicircum{}λ\_j=-∂\_iV̇\_j, similarly -∂\_jV̇\_i. Last two vanish. Hence R\_\{ij\}=-(1/2)(∂\_iV̇\_j+∂\_jV̇\_i). ✓

  Type: claim, Inference: assumption, Status: validated, Taint: clean

\item \textbf{Node 1.4.4:} 4.1.4 (Vector R): R=∂\_μ∂\_νh\textasciicircum{}\{μν\}-□h. h\textasciicircum{}\{00\}=0, h\textasciicircum{}\{0i\}=-V\_i, h\textasciicircum{}\{ij\}=0. ∂\_μ∂\_νh\textasciicircum{}\{μν\}=-2∂\_i∂\_0V\_i=0 (transversality). □h=0. Hence R=0. ✓

  Type: claim, Inference: assumption, Status: validated, Taint: clean

\item \textbf{Node 1.4.5:} 4.2.1-4.2.4 (Tensor sector): Set h\_\{ij\}=h\textasciicircum{}\{TT\}\_\{ij\}, h\_\{00\}=h\_\{0i\}=0. R\_\{00\}: h\textasciicircum{}0\_0=h\textasciicircum{}i\_0=0, both terms vanish. R\_\{0i\}: ∂\_j∂\_0h\textasciicircum{}\{TT\}\_\{ji\}=∂\_0(∂\_jh\textasciicircum{}\{TT\}\_\{ji\})=0 by transversality. R\_\{ij\}: 2R\_\{ij\}=∂\_k∂\_ih\textasciicircum{}\{TT\}\_\{kj\}+∂\_k∂\_jh\textasciicircum{}\{TT\}\_\{ki\}-∂\_i∂\_jh\textasciicircum{}\{TT\}\_\{kk\}-□h\textasciicircum{}\{TT\}\_\{ij\}. First two: transversality→0. Third: tracelessness→0. Hence R\_\{ij\}=-(1/2)□h\textasciicircum{}\{TT\}\_\{ij\}. R: ∂\_i∂\_jh\textasciicircum{}\{TT\}\_\{ij\}=0 (transversality²), □(h\textasciicircum{}\{TT\}\_\{ii\})=0 (tracelessness). ✓ QED.

  Type: claim, Inference: assumption, Status: validated, Taint: clean

\end{enumerate}
\item \textbf{Node 1.5:} Theorem 5.1: The action splits as I = I\_\{TT\} + I\_V + I\_S with no cross-terms: I\_\{TT\} = (1/4)∫(□h\textasciicircum{}\{TT\}\_\{ij\})², I\_V = -(1/2)∫V\_i∇²□V\_i, I\_S = ∫(Φ*,ψ*)M(Φ,ψ)d⁴k/(2π)⁴ with M the 2×2 matrix. No cross-terms by Schur orthogonality for SO(3) irreps (spin-0, spin-1, spin-2). [Refs: Schur Orthogonality for SO(3) Irreps. Depends: 1.3, 1.4]

  Type: claim, Inference: assumption, Status: validated, Taint: clean

\begin{enumerate}
\item \textbf{Node 1.5.1:} 5.1.1 (No cross-terms): The scalar, vector, and tensor sectors transform as spin-0, spin-1, spin-2 under SO(3). Any quadratic form Q[h]=∫h\_\{μν\}O\textasciicircum{}\{μνρσ\}h\_\{ρσ\} with O rotation-covariant decomposes with no cross-terms by Schur orthogonality. [Ref: Schur Orthogonality for SO(3) Irreps]

  Type: claim, Inference: assumption, Status: validated, Taint: clean

\item \textbf{Node 1.5.2:} 5.1.2 (Tensor sector): R\_\{μν\}R\textasciicircum{}\{μν\}|\_\{TT\}=R\_\{ij\}R\_\{ij\}|\_\{TT\} (since R\_\{00\}, R\_\{0i\} vanish by 4.2). By Claim 4.2: =(1/4)(□h\textasciicircum{}\{TT\}\_\{ij\})². The β-term vanishes since R|\_\{TT\}=0. Hence I\_\{TT\}=(1/4)∫(□h\textasciicircum{}\{TT\}\_\{ij\})². ✓

  Type: claim, Inference: assumption, Status: validated, Taint: clean

\item \textbf{Node 1.5.3:} 5.1.3 (Vector sector): R\_\{μν\}R\textasciicircum{}\{μν\}=(R\_\{00\})²-2R\_\{0i\}R\_\{0i\}+R\_\{ij\}R\_\{ij\} with signature (-,+,+,+). R\_\{00\}²=0; R\_\{0i\}R\_\{0i\}=(1/4)(∇²V\_i)²; ∫R\_\{ij\}R\_\{ij\}=(1/2)∫(∂\_iV̇\_j)² (cross-term vanishes by transversality). After IBP: ∫R\_\{μν\}R\textasciicircum{}\{μν\}|\_V=-(1/2)∫V\_i∇²□V\_i. β-term vanishes since R|\_V=0. ✓

  Type: claim, Inference: assumption, Status: validated, Taint: clean

\item \textbf{Node 1.5.4:} 5.1.4 (Scalar sector): Substitute curvature components from Claim 3.1 into R\_\{μν\}R\textasciicircum{}\{μν\}-βR². Pass to Fourier: ∇²→-k², □→-p². Collect all terms in |Φ|², Re(Φ*ψ), |ψ|² to obtain the matrix M. [Verified by SymPy test\_scalar\_action.py] ✓ QED.

  Type: claim, Inference: assumption, Status: validated, Taint: clean

\end{enumerate}
\item \textbf{Node 1.6:} Lemma 6.1 + Theorem 6.3: det(M) = 8(1-3β)k⁴p⁴. For β≠1/3, the scalar propagators are G = (1/2)M⁻¹ giving ⟨ΦΦ⟩ = 3/(4k⁴) + 1/(2k²p²) + (1-2β)/(8(1-3β)p⁴), ⟨ψψ⟩ = (1-2β)/(8(1-3β)p⁴), ⟨Φψ⟩ = -1/(4k²p²) - (1-2β)/(8(1-3β)p⁴). [Algebraic computation, verified by SymPy. Depends: 1.5]

  Type: claim, Inference: assumption, Status: validated, Taint: clean

\begin{enumerate}
\item \textbf{Node 1.6.1:} 6.1.1-6.1.3: Write M\_\{11\}=2(1-2β)k⁴, M\_\{12\}=4(1-3β)k²p²+2(1-2β)k⁴, M\_\{22\}=12(1-3β)p⁴+8(1-3β)k²p²+2(1-2β)k⁴. Expand M\_\{11\}M\_\{22\} and M\_\{12\}². [Algebraic computation]

  Type: claim, Inference: assumption, Status: validated, Taint: clean

\item \textbf{Node 1.6.2:} 6.1.4-6.1.5: Subtracting M\_\{11\}M\_\{22\}-M\_\{12\}²: k⁸ terms cancel; k⁶p² terms cancel (both 16(1-2β)(1-3β)); k⁴p⁴ terms give [24(1-2β)(1-3β)-16(1-3β)²]k⁴p⁴ = 8(1-3β)[3(1-2β)-2(1-3β)]k⁴p⁴ = 8(1-3β)·1·k⁴p⁴. Hence det(M)=8(1-3β)k⁴p⁴. [Verified by SymPy test\_determinant.py] ✓

  Type: claim, Inference: assumption, Status: validated, Taint: clean

\item \textbf{Node 1.6.3:} 6.3.1 (Normalization): Action I\_S=∫Φ*\_a M\_\{ab\} Φ\_b over all k. For real fields Φ\_a(-k)=Φ*\_a(k), so canonical form is (1/2)∫Φ(-k)A(k)Φ(k) with A=2M. Propagator G=A⁻¹=(1/2)M⁻¹. Consistency: tensor gives G=2/p⁴ ✓, vector gives G=1/(k²p²) ✓.

  Type: claim, Inference: assumption, Status: validated, Taint: clean

\item \textbf{Node 1.6.4:} 6.3.4-6.3.7: (1/2)M⁻¹ = (1/(2detM))[[M\_\{22\},-M\_\{12\}],[-M\_\{12\},M\_\{11\}]]. ⟨ψψ⟩=M\_\{11\}/(2detM)=(1-2β)/(8(1-3β)p⁴). ⟨Φψ⟩=-M\_\{12\}/(2detM)=-1/(4k²p²)-(1-2β)/(8(1-3β)p⁴). ⟨ΦΦ⟩=M\_\{22\}/(2detM)=3/(4k⁴)+1/(2k²p²)+(1-2β)/(8(1-3β)p⁴). [Verified by SymPy test\_matrix\_inverse.py] ✓ QED.

  Type: claim, Inference: assumption, Status: validated, Taint: clean

\end{enumerate}
\item \textbf{Node 1.7:} Theorems 6.4-6.5 + Remarks 7.1-7.2: Vector propagator ⟨V\_iV\_j⟩ = P\textasciicircum{}T\_\{ij\}/(k²p²). Tensor propagator ⟨h\textasciicircum{}\{TT\}\_\{ij\}h\textasciicircum{}\{TT\}\_\{kl\}⟩ = 2Π\textasciicircum{}\{TT\}\_\{ijkl\}/p⁴. The 1/p⁴ poles are the hallmark of fourth-derivative gravity (massless graviton + Weyl ghost). At β=1/3, det(M)=0 (conformal gravity). At β=1/2, the scalar sector reduces to one constrained DOF. [Depends: 1.5, 1.6]

  Type: claim, Inference: assumption, Status: validated, Taint: clean

\begin{enumerate}
\item \textbf{Node 1.7.1:} 6.4.1-6.4.2 (Vector propagator): In Fourier, I\_V=(1/2)∫k²p²V*\_i V\_i restricted to transverse modes (k\_iV\_i=0). The operator A=k²p². The inverse on the transverse subspace: A⁻¹P\textasciicircum{}T\_\{ij\}=P\textasciicircum{}T\_\{ij\}/(k²p²). Hence ⟨V\_iV\_j⟩=P\textasciicircum{}T\_\{ij\}/(k²p²). ✓

  Type: claim, Inference: assumption, Status: validated, Taint: clean

\item \textbf{Node 1.7.2:} 6.5.1-6.5.2 (Tensor propagator): In Fourier, I\_\{TT\}=(1/4)∫p⁴h\textasciicircum{}\{TT*\}\_\{ij\}h\textasciicircum{}\{TT\}\_\{ij\}. Identifying A/2=p⁴/4 gives A=p⁴/2, so G=2/p⁴. The TT projector Π\textasciicircum{}\{TT\}\_\{ijkl\} is the identity on TT symmetric tensors. Hence ⟨h\textasciicircum{}\{TT\}\_\{ij\}h\textasciicircum{}\{TT\}\_\{kl\}⟩=2Π\textasciicircum{}\{TT\}\_\{ijkl\}/p⁴. ✓

  Type: claim, Inference: assumption, Status: validated, Taint: clean

\item \textbf{Node 1.7.3:} Remark 7.1: The 1/p⁴ poles correspond to a double pole at p²=0, splitting into a massless graviton and a massless Weyl ghost in the Hamiltonian decomposition. The 1/k⁴ piece in ⟨ΦΦ⟩ is instantaneous (no ω-dependence) — the fourth-derivative analogue of the Newtonian potential constraint.

  Type: claim, Inference: assumption, Status: validated, Taint: clean

\item \textbf{Node 1.7.4:} Remark 7.2: At β=1/2: (1-2β)=0, so ⟨ψψ⟩=0, p⁻⁴ pieces in ⟨ΦΦ⟩ and ⟨Φψ⟩ vanish. ⟨Φψ⟩=-1/(4k²p²), ⟨ΦΦ⟩=3/(4k⁴)+1/(2k²p²). At β=1/3: det(M)=0, the conformal gravity point where the Weyl tensor squared action has conformal symmetry removing one scalar DOF. [Ref: Gauss-Bonnet at Linear Order. Verified by SymPy test\_special\_cases.py] ✓ QED.

  Type: claim, Inference: assumption, Status: validated, Taint: clean

\end{enumerate}
\end{enumerate}



\end{document}
